\clearpage
\VerdierPhysicalPageAnchor{0171}%
\begingroup\small\setlength{\parskip}{0pt}\setlength{\emergencystretch}{1.5em}
\VerdierSemanticAnchor{ch4-4-3-3}
\medskip\noindent\textbf{4.3.3.}\quad
On s'intéresse plus particulièrement aux termes
$E\!\left[\begin{smallmatrix}n\\p-r+1,p,p+1,p+r\end{smallmatrix}\right]$,
$p\in\mathbb Z$, $r\in\mathbb Z$, $r>0$, car ceux-ci permettront, sous des
hypothèses convenables, de déterminer le gradué associé à
$H^n(-\infty,+\infty)$. Pour nous ramener aux notations de \emph{loc. cit.},
nous poserons :

\VerdierSemanticAnchor{ch4-4-3-3-1}
\[
\left\{
\begin{aligned}
E\!\left[\begin{smallmatrix}n\\p-r,p-1,p,p+r-1\end{smallmatrix}\right]
  &=\mathrm I_r^{-p,n+p},\\
&\textit{i.e.}:\\[-2pt]
\mathrm I_r^{p,q}
  &=E\!\left[\begin{smallmatrix}p+q\\-p-r,-p-1,-p,-p+r-1\end{smallmatrix}\right],
  &&p,q\in\mathbb Z,\quad 1\leq r\leq+\infty,\\
&\text{Les différentielles correspondantes seront notées :}\\[-2pt]
d_r\colon\mathrm I_r^{p,q}&\longrightarrow
  \mathrm I_r^{p+r,q-r+1}.
\end{aligned}
\right.
\tag{4.3.3.1}
\]

Il nous arrivera très souvent d'utiliser la notation :

\VerdierSemanticAnchor{ch4-4-3-3-2}
\[
\left\{
\begin{aligned}
E\!\left[\begin{smallmatrix}n\\p-r+1,p,p+1,p+r\end{smallmatrix}\right]
  &=\mathrm{II}_{r+1}^{n-p,p},\\
&\textit{i.e.}:\\[-2pt]
\mathrm{II}_r^{p,q}
  &=E\!\left[\begin{smallmatrix}p+q\\q-r+2,q,q+1,q+r-1\end{smallmatrix}\right],
  &&p,q\in\mathbb Z,\quad 2\leq r\leq+\infty,\\
&\text{Les différentielles correspondantes seront notées :}\\[-2pt]
d_r\colon\mathrm{II}_r^{p,q}&\longrightarrow
  \mathrm{II}_r^{p+r,q-r+1}.
\end{aligned}
\right.
\tag{4.3.3.2}
\]

Nous préciserons dans chaque cas la notation utilisée. Remarquons que
lorsqu'on utilise la notation (4.3.3.1), les termes $\mathrm I_r^{p,q}$ sont
définis pour $r\geq1$, alors que lorsqu'on utilise la notation (4.3.3.2),
les termes $\mathrm{II}_r^{p,q}$ ne sont définis que pour $r\geq2$.

\VerdierSemanticAnchor{ch4-4-3-3-spectral-sequence}
L'ensemble des termes
$E\!\left[\begin{smallmatrix}n\\p-r+1,p,p+1,p+r\end{smallmatrix}\right]$,
$n,p\in\mathbb Z$, $1\leq r\leq+\infty$, munis des différentielles
correspondantes, est appelé la \emph{suite spectrale associée à l'objet
spectral} $(H^n(p,q),\delta(p,q,r))$.
\endgroup
